\lecture{4}{Fri 09 Sep 2022 10:29}{Consequences of Lagrange's Theorem}

\begin{corollary}
	(Theorem 2.4.3) Any group of prime order has no proper subgroups, and in particular is cyclic.
\end{corollary}
\begin{proof}
	Say $|G| = p$ is some prime number, suppose $H\le G$. Then $|H| \mid p$, so $|H| = 1$ or $p$, whence $H$ must be $G$ or $\{e\} $.

	But if $H = \left\langle a \right\rangle $ for some nonidentity $a$, then $H \le G$. We have $H \neq \{e\} $. Hence $H = G$. Thus $G$ is cyclic.
\end{proof}
\begin{remark}
	This means groups of prime order are \textbf{abelian} as well.
\end{remark}

\begin{definition}[Order of an element]
	Let $G$ be a group. The smallest positive integer $n$ such that $a^n = e$, if one exists, is called the \textit{order} of $a \in G$, and is written $o(a)$. If no such integer exists, we define $o(a) = +\infty $.
\end{definition}

\begin{example}
	\begin{enumerate}
		\item In $\left( \mathbb{R}^\times  \right) $, one has $o(-1) = 2$, because $(-1)^{2} = 1$.
		\item In $\mathbb{Z}$, $o(1) = +\infty$.
		\item In $\mathbb{Z}_3$, $o(1) = 3$ since $1+1+1 \equiv_3 0$.
	\end{enumerate}
\end{example}

\begin{theorem}
	When $G$ is a finite group, we have $o(a) \mid |G|$ for all $a \in G$.
\end{theorem}
\begin{proof}
If $a \in G$, we have $o(a) = |\left\langle a \right\rangle |$. Thus, by Lagrange's Theorem, we have $o(a) \mid |G|$ since $|\left\langle a \right\rangle | \mid |G|$.
\end{proof}

\begin{corollary}
	Suppose that  $G$ is a finite group. Then for all $a \in G$, \[
		a^{|G|} = e
	.\] 
\end{corollary}
\begin{proof}
	Since $o(a) \mid |G|$, we have $|G| = m o(a)$. Now \[
		a^{|G|} = \left( a^{o(a)} \right) ^{m} = e^{m} = e
	.\] 
\end{proof}

\subsubsection{Elementary Number Theory}

We've seen $\left( \mathbb{Z}_n, + \right) $. We think of it a equivalence class of integers: \[
	\mathbb{Z}_n = \{[0],[1],\ldots,[n-1]\} 
\] 
where the equivalence relation is integers modulo $n$.

$[a]+[b] = [a+b]$

We can also consider multiplication (mod n).
 \[
	 [a] \times [b] = [ab]
.\] 

\begin{itemize}
	\item (identity) $[1]$ is identity.
	\item (inverse) $[2]^{-1}$ does not exist in $\mathbb{Z}_6$. 
\end{itemize}

Define
\[
	U_n = \{[a]:(a,b) = 1\} 
.\]
\begin{notation}
	In number theory we often write $U_n$ as  $\left( \mathbb{Z}/ n\mathbb{Z} \right) ^{\times }$
\end{notation}
\begin{claim}
	$U_n$ forms a group with multiplication mod $n$.
\end{claim}
\begin{proof}
	\begin{itemize}
		\item (closure)
	If $(a,n) = 1$ and $(b,n) = 1$ then $(ab, n) = 1$. Hence $[a][b] = [ab] \in U_n$.
\item (identity) $[1]$ is identity.
\item (associativity) Inherited from $\left( \mathbb{Q}^{x}, \times  \right) $
\item (inverse) If $(a,n) = 1$, then there exists integers $x, y \in \mathbb{Z}$ such that $ax + ny = 1$ (Euclidean Algorithm), whence $ax\equiv_n 1$.

	Then $[x] = [a]^{-1}$.
	So $\left( U_n, \times  \right) $ is indeed a group.
	\end{itemize}
\end{proof}

Hence if $a \in U_n$, then $a^{|U_n|} \equiv_n 1$. 

\begin{definition}[Euler  $\varphi$-funtion]
The Euler $\varphi$-function is defined for each $n \in \mathbb{N}$ by \[
	\varphi(n) = \left| \{1\le a\le n: (a,n) = 1\}  \right| 
.\] 
\end{definition}
\begin{fact}
	\[
		\phi(n) = n \prod_{p^4 \mid n, p \text{prime}} (1-\frac{1}{p})  
	\]
\end{fact}

\begin{corollary}[Euler's Theorem]
	When $(a,n) = 1$ one has $a^{\phi(n)} \equiv_n 1$.
\end{corollary}

\begin{corollary}[Fermat's Little Theorem]
	When $p$ is prime, one has $a^p \equiv_p a$ for all  $a \in \mathbb{Z}$.
\end{corollary}
\begin{proof}
	Apply Euler's Theorem when $(a,p) = 1$.
	When $(a,p) = 1$, If $p|a$, then $a^p \equiv a \equiv 0$ (mod $p$ ).
\end{proof}

Which $U_n$ are cyclic?

Onserve: $U_n$ is abelian. 

\begin{fact}
	Every finite abelian group is "isomorphic" to a group of the shape \[
		U_n^{(d)} = \{a^d: a \in U_n\} 
	\] 
	for some $d, n \in \mathbb{N}$.
\end{fact}

\subsection{Homomorphisms and normal subgroups}

\begin{definition}[Homomorphism]
	Let $G$ and $G'$ be two groups. The mapping \[
		\phi: G \to  G'
	\] 
	is called a \textit{homomorphism} when for all $a,b \in G$, one has $\phi(ab) = \phi(a)\phi(b)$.

	\begin{itemize}
		\item 
	Injective homomorphism is called \textit{monomorphism}. 
\item Surjective homomorphism is called \textit{epimorphism}. 
\item Bijective homomorphism is called \textit{isomorphism}.
\item Bijective self-mapping homomorphism is called an \textit{automorphism}.
	\end{itemize}
\end{definition}

\begin{definition}
	We say that two groups $G$ and $G'$ are \textit{isomorphic} if there exists a mapping $\phi: G \to G'$ such that $\phi$ is an isomorphism. We write \[
		G \cong G'
	.\] 
\end{definition}

